\chapter{Kohomologia z nośnikami zwartymi i homologia Borel--Moore'a}
\label{ch:compact-bm-v11}

\section{Motywacja: dlaczego zwykła dualność nie wystarcza?}
Dla zwartej zorientowanej $n$-rozmaitości $M$ znamy izomorfizm
$H^k(M;\Z)\cong H_{n-k}(M;\Z)$. Na prostej $\R$ nie można
bez zmian napisać tego samego: $H^0(\R;\Z)=\Z$, lecz
$H_1(\R;\Z)=0$. Intuicja jest geometryczna.
Otwarta prosta ma klasę fundamentalną reprezentowaną przez
\emph{nieskończoną, ale lokalnie skończoną} sumę kolejnych
zorientowanych odcinków. Nie stanowi ona zwykłego skończonego
łańcucha singularnego. Dualnie potrzebujemy form i kołańcuchów,
które mają nośnik zwarty i mogą być całkowane po takiej sumie.
Otrzymamy dwa różne izomorfizmy, których nie wolno mylić:
\[
H_c^k(M;G)\cong H_{n-k}(M;G),\qquad
H^k(M;G)\cong H_{n-k}^{\mathrm{BM}}(M;G)
\]
(dla rozmaitości zorientowanych, bez brzegu; w przypadku
nieorientowalnym należy skręcić współczynniki systemem orientacji).

\section{Kohomologia z nośnikami zwartymi}
\begin{definicja}[Definicja przez pary względne]
\label{def:compact-support-cohom-v11}
Niech $X$ będzie lokalnie zwartą przestrzenią Hausdorffa,
a $G$ grupą abelową. Przez $\mathcal K(X)$ oznaczmy
skierowany zbiór zwartych podzbiorów $X$ uporządkowany
przez inkluzję. Definiujemy
\[
H_c^k(X;G):=\varinjlim_{K\in\mathcal K(X)}
H^k(X,X\setminus K;G).
\]
Jeżeli $K\subseteq L$, to zawieranie par
$(X,X\setminus L)\to(X,X\setminus K)$ daje przez
kontrawariantność kohomologii morfizm
$H^k(X,X\setminus K;G)\to H^k(X,X\setminus L;G)$.
To właśnie są strzałki układu skierowanego.
\end{definicja}

\begin{uwaga}[Nośnik kołańcucha a nośnik klasy]
Klasa $H_c^k$ nie oznacza dowolnego kołańcucha singularnego
zerowego poza skończoną listą sympleksów. Definicja względna
mówi, że klasa może być reprezentowana \emph{względem dopełnienia}
pewnego zwartego $K$. Dla lokalnie skończonych triangulacji
rozmaitości porównamy ją z kołańcuchami komórkowymi o skończonym
nośniku. W rozdziale o kohomologii snopów (Rozdział~\ref{ch:snopy-verdier-v12}) rozróżniamy
ściśle: nośnik sekcji snopu, nośnik względnej klasy
singularnej oraz lokalnie skończony nośnik łańcucha Borel--Moore'a. Ta równoważność wymaga twierdzenia o porównaniu
konstrukcji komórkowej i singularnej.
\end{uwaga}

\begin{stwierdzenie}[Przestrzeń zwarta i mapa właściwa]
\label{prop:compact-proper-v11}
Jeśli $X$ jest zwarta, to $H_c^k(X;G)\cong H^k(X;G)$.
Jeżeli $f:X\to Y$ jest ciągłym odwzorowaniem \emph{właściwym},
tj. przeciwobraz każdego zwartego podzbioru $Y$ jest zwarty,
to $f$ indukuje homomorfizm
$f^*:H_c^k(Y;G)\to H_c^k(X;G)$.
\end{stwierdzenie}
\begin{proof}
Dla zwartego $X$ zbiór $K=X$ jest największym elementem
$\mathcal K(X)$; granica prosta równa się
$H^k(X,X\setminus X;G)=H^k(X;G)$.
Dla mapy właściwej i zwartego $K\subseteq Y$ mamy
$L=f^{-1}(K)$ zwarty, a mapę par
$f:(X,X\setminus L)\to(Y,Y\setminus K)$.
Daje ona morfizm w kohomologii względnej. Naturalność
względem inkluzji zbiorów zwartych pozwala przejść
do granic prostych. Bez właściwości $f$ przeciwobraz
$K$ może nie być zwarty, więc konstrukcja nie musi działać.
\end{proof}

\begin{przyklad}[Dlaczego $H_c^0(\R)=0$?]
Na spójnej niezwartej rozmaitości bez brzegu
każda lokalnie stała funkcja o nośniku zwartym jest zerowa.
Z tego powodu $H_c^0(\R;\Z)=0$, podczas gdy
$H^0(\R;\Z)=\Z$. Na przestrzeni dyskretnej $\N$
klasa w stopniu zerowym może być niezerowa tylko na
skończenie wielu punktach, zatem
$H_c^0(\N;\Z)\cong\bigoplus_{m\in\N}\Z$.
\end{przyklad}

\section{Kompaktowo podparte formy różniczkowe}
\begin{definicja}[Kompleks de Rhama z nośnikami zwartymi]
Dla gładkiej rozmaitości $M$ niech $\Omega_c^k(M)$ będzie
przestrzenią gładkich $k$-form $\omega$ o zwartym nośniku
$\operatorname{supp}\omega=\overline{\{x:\omega_x\ne0\}}$.
Ponieważ $\operatorname{supp}(\dd\omega)\subseteq
\operatorname{supp}(\omega)$, różniczka zewnętrzna definiuje
kompleks $(\Omega_c^*(M),\dd)$. Jego kohomologię oznaczamy
$H_{c,\mathrm{dR}}^k(M)$.
\end{definicja}

\begin{lemat}[Całkowanie na prostej]
\label{lem:compact-R-v11}
Odwzorowanie $\omega=f(x)\,\dd x\mapsto
\int_{x\in\R}f(x)\,\dd x$ indukuje izomorfizmy
$H_{c,\mathrm{dR}}^1(\R)\cong\R$ oraz
$H_{c,\mathrm{dR}}^0(\R)=0$.
\end{lemat}
\begin{proof}
Jeśli $f\in C_c^\infty(\R)$, to $f\,\dd x$ jest zamknięta.
Jeżeli $f\,\dd x=\dd F$ dla pewnej
$F\in C_c^\infty(\R)$, to
$\int_{x\in\R}f(x)\,\dd x=F(+\infty)-F(-\infty)=0$.
Odwrotnie, gdy ta całka jest zerowa, funkcja
$F(x)=\int_{u\in(-\infty,x]} f(u)\,\dd u$ ma nośnik
zwarty: na lewo od nośnika $f$ i na prawo od niego
jest równa zeru. Zatem $\dd F=f\,\dd x$.
Istnieje $f$ o całce $1$, co dowodzi surjektywności.
Dla form stopnia $0$ warunek $\dd F=0$ oznacza
stałość $F$, a jedyna stała funkcja o nośniku zwartym
na $\R$ jest zerowa.
\end{proof}

\begin{przyklad}[Właściwość odwzorowania i stopień na prostej]
\label{prz:compact-proper-degree-22}
Na $\R$ zorientowanej standardowo wybierzmy
$\rho\in C_c^\infty(\R)$ z
$\int_{y\in\R}\rho(y)\,\dd y=1$.
Lemat~\ref{lem:compact-R-v11} mówi, że
$[\rho(y)\,\dd y]$ generuje
$H_{c,\mathrm{dR}}^1(\R)\cong\R$.
Rozpatrzmy dwa odwzorowania właściwe:
$f_2(x)=x^2$ i $f_3(x)=x^3$.
Właściwość wynika z tego, że przeciwobraz
każdego zbioru zwartego jest domknięty i ograniczony.
Działanie na reprezentancie generatora daje
\[
 f_2^*(\rho(y)\,\dd y)=2x\rho(x^2)\,\dd x,
 \qquad
 f_3^*(\rho(y)\,\dd y)=3x^2\rho(x^3)\,\dd x.
\]
Pierwszy współczynnik jest nieparzysty,
więc jego całka po $\R$ znika. Dla drugiej
formy podstawienie $y=x^3$ daje całkę równą $1$.
Stąd $(f_2)^*=0$, a $(f_3)^*=\id$ na
$H_{c,\mathrm{dR}}^1(\R)$. Zatem odwzorowanie
właściwe może działać zerowo na kohomologii
z nośnikami zwartymi; jego działanie odzwierciedla
stopień, a nie samą właściwość.

Dla kontrastu włączenie
$j:(0,1)\hookrightarrow\R$ nie jest właściwe:
przeciwobraz $[0,1]$ to $(0,1)$, który nie jest
zwarty. Wybierając $\rho$ niezerową w otoczeniu
$0$, otrzymujemy formę $j^*(\rho(y)\,\dd y)$,
której nośnik w $(0,1)$ nie jest zwarty.
Dlatego ogólne odwzorowanie niewłaściwe nie definiuje
przeciwobrazu na całym kompleksie form o nośniku zwartym.
\end{przyklad}

\begin{twierdzenie}[Lokalny rachunek na $\R^n$]
\label{thm:compact-Rn-v11}
Dla $n\ge1$ zachodzi
$H_{c,\mathrm{dR}}^k(\R^n)=0$ przy $k\ne n$ oraz
$H_{c,\mathrm{dR}}^n(\R^n)\cong\R$.
W stopniu $n$ izomorfizm jest dany przez całkowanie form.
\end{twierdzenie}
\begin{proof}
Przedstawimy indukcję wykorzystującą $\R^n=\R^{n-1}\times\R$.
Wybierzmy $\rho\in C_c^\infty(\R)$ o
$\int_{t\in\R}\rho(t)\,\dd t=1$.
Rozłóżmy formę w sposób jednoznaczny:
$\omega=\alpha(x,t)+\dd t\wedge\beta(x,t)$,
gdzie $\alpha$ i $\beta$ nie zawierają $\dd t$.
Operator całkowania po włóknie określamy wzorem
$I(\omega)=\int_{t\in\R}\beta(x,t)\,\dd t$.
Jest to forma stopnia o jeden mniejszego, o nośniku zwartym,
na $\R^{n-1}$. Zapisując $\dd_x$ dla różniczki względem $x$,
mamy
$\dd\omega=\dd_x\alpha+\dd t\wedge
(\partial_t\alpha-\dd_x\beta)$, a więc
$I\dd=-\dd_x I$. Definiujemy też
$s(\eta)=\rho(t)\,\dd t\wedge\eta$;
wówczas $I\circ s=\id$ i $\dd s=-s\dd_x$.

Dla formy poziomej $\beta$ niech
$P\beta(x,t)=\int_{u\in(-\infty,t]}\beta(x,u)\,\dd u-
\bigl(\int_{u\in\R}\beta(x,u)\,\dd u\bigr)R(t)$,
gdzie $R(t)=\int_{u\in(-\infty,t]}\rho(u)\,\dd u$.
Wynik ma nośnik zwarty, gdyż poza dużym przedziałem
oba wyrazy się znoszą, również jednostajnie względem $x$.
Określmy $H(\omega)=P\beta$. Mamy
$\partial_t P\beta=\beta-\rho\,I\omega$ i
$P(\partial_t\alpha)=\alpha$ dla form $\alpha$ o
nośniku zwartym, a $P\dd_x=\dd_xP$.
Stąd przez bezpośrednie podstawienie
$\dd H\omega+H\dd\omega=
\alpha+\dd t\wedge(\beta-\rho\,I\omega)
=\omega-sI\omega$.
Jest to jawna homotopia łańcuchowa; mimo że $I$ i $s$
zmieniają stopień o jeden, indukują odwrotne
izomorfizmy kohomologii po przesunięciu stopnia. Zatem
$H_{c,\mathrm{dR}}^k(\R^n)\cong
H_{c,\mathrm{dR}}^{k-1}(\R^{n-1})$.
Początek indukcji daje Lemat~\ref{lem:compact-R-v11}.
\end{proof}

\begin{twierdzenie}[De Rham z nośnikami zwartymi]
\label{thm:compact-deRham-v11}
Dla gładkiej drugoprzeliczalnej rozmaitości $M$ całkowanie
form po gładkich sympleksach singularnych wyznacza
naturalne izomorfizmy
$H_{c,\mathrm{dR}}^k(M)\cong H_c^k(M;\R)$.
\end{twierdzenie}
\begin{proof}[Argument wykorzystujący globalne twierdzenie porównawcze]
Wyjaśniamy zakres porównania. Na otwartym zbiorze
współrzędnych izomorfizm otrzymuje się ze zwykłego
lematu Poincarégo i wersji z nośnikami zwartymi
z Twierdzenia~\ref{thm:compact-Rn-v11}.
Dla otwartego pokrycia $M=U\cup V$ istnieje krótki ciąg
kompleksów
$0\to\Omega_c^*(U\cap V)\to
\Omega_c^*(U)\oplus\Omega_c^*(V)\to
\Omega_c^*(M)\to0$:
pierwsza mapa to $(\alpha,-\alpha)$ przez przedłużenie
zerem, druga jest sumą, a surjektywność wynika z podziału
jedności na zwartym nośniku formy. Prowadzi to do ciągu
Mayera--Vietorisa. Analogiczny długi ciąg Mayera--Vietorisa istnieje
w singularnej kohomologii z nośnikami zwartymi,
ale nie wynika z samej definicji przez pary:
trzeba najpierw zastosować ekscyzję i odpowiednią
konstrukcję z nośnikami (lub rezolwenty snopowe).
Porównanie map łączących i lemat pięciu przenoszą
izomorfizm z otwartych kostek na ich skończone sumy,
a następnie na wszystkie rozmaitości przy użyciu
lokalnie skończonego pokrycia i wyczerpania przez
zwarte podzbiory. Ostatni krok można sformułować jako
standardowe twierdzenie o rozdzielczości kompleksu
stałych snopów przez formy różniczkowe; nie utożsamiamy
go z pojedynczym zastosowaniem lematu Poincarégo.
\end{proof}

\section{Łańcuchy lokalnie skończone i homologia Borel--Moore'a}
\begin{definicja}[Kompleks lokalnie skończonych łańcuchów]
\label{def:BM-chain-v11}
Niech $X$ będzie lokalnie zwartą przestrzenią Hausdorffa.
\emph{Lokalnie skończonym singularnym $k$-łańcuchem}
o współczynnikach w $G$ nazywamy formalną sumę
$c=\sum_{\sigma:\Delta^k\to X}a_\sigma\sigma$,
w której dla każdego zwartego $K\subset X$ tylko
skończenie wiele sympleksów o $a_\sigma\ne0$
ma obraz przecinający $K$. Zbiór takich łańcuchów
oznaczamy $C_k^{\mathrm{lf}}(X;G)$.
Różniczkę brzegową definiujemy na składnikach jak
w homologii singularnej; jest dobrze określona,
ponieważ obraz ściany sympleksu zawiera się w obrazie
samego sympleksu. Homologię tego kompleksu oznaczamy
$H_k^{\mathrm{BM}}(X;G)$ i nazywamy
\emph{homologią Borel--Moore'a} (lub homologią lokalnie skończoną).
\end{definicja}

\begin{stwierdzenie}[Zwartość i właściwość]
\label{prop:BM-proper-v11}
Jeżeli $X$ jest zwarta, to $H_k^{\mathrm{BM}}(X;G)=H_k(X;G)$.
Jeżeli $f:X\to Y$ jest właściwym odwzorowaniem ciągłym
między lokalnie zwartymi przestrzeniami Hausdorffa,
to indukuje $f_*:H_k^{\mathrm{BM}}(X;G)\to H_k^{\mathrm{BM}}(Y;G)$.
\end{stwierdzenie}
\begin{proof}
W przestrzeni zwartej warunek lokalnej skończoności
zastosowany do $K=X$ wymusza skończoność całej sumy.
Jeśli $c$ jest lokalnie skończony, a $L\subset Y$ jest zwarty,
to $f^{-1}(L)$ jest zwarty. Tylko skończenie wiele
sympleksów z $c$ przecina ten przeciwobraz, toteż
ich obrazy przez $f$ przecinają $L$ jedynie skończenie
wiele razy; po zsumowaniu współczynników jednakowych
obrazów nadal otrzymujemy rodzinę lokalnie skończoną.
Mapa łańcuchowa jest zatem określona
poprawnie. Przy niewłaściwych mapach ta własność
może zawieść; np. mapa stała $\R\to\{*\}$
wysyłałaby nieskończony lokalnie skończony łańcuch
na nieskończoną sumę w punkcie.
\end{proof}

\begin{center}
\begin{tikzpicture}[>=Stealth,font=\small,x=0.8cm,y=0.8cm]
\fill[black!7] (-1.05,-.55) rectangle (1.05,.50);
\draw[thick] (-4.4,0)--(4.4,0);
\foreach \m in {-4,-3,-2,-1,0,1,2,3,4}
 {\fill[manifoldblue] (\m,0) circle (1.5pt);}
\foreach \m in {-4,-3,-2,-1,0,1,2,3}
 {\draw[->,manifoldblue,semithick] (\m+.12,.12)--(\m+.86,.12);}
\node[anchor=east] at (-4.5,.09) {$\cdots$};
\node[anchor=west] at (4.5,.09) {$\cdots$};
\node[below] at (-1,-.10) {$-1$};
\node[below] at (0,-.10) {$0$};
\node[below] at (1,-.10) {$1$};
\node[above] at (0,.53) {$K=[-1,1]$};
\end{tikzpicture}
\par\smallskip\small Schemat lokalnie skończonego łańcucha na prostej:
każdy zwarty odcinek przecina tylko skończenie wiele składników.
\end{center}

\begin{przyklad}[Nieskończony cykl na prostej]
\label{ex:BM-line-v11}
Niech $\sigma_m:[0,1]\to\R$ będzie dane przez
$\sigma_m(t)=m+t$, $m\in\Z$. Suma
$c=\sum_{m\in\Z}\sigma_m$ jest lokalnie skończona;
jej brzeg znika, ponieważ przy każdym wierzchołku
$m\in\Z$ wkład końca $\sigma_{m-1}$ znosi wkład
początku $\sigma_m$. Nie jest to skończony łańcuch
singularny. Reprezentuje generator
$H_1^{\mathrm{BM}}(\R;\Z)\cong\Z$.
Natomiast łańcuch $\sum_{m\ge0}\sigma_m$ ma brzeg
$-[0]$, więc $[0]$ znika w $H_0^{\mathrm{BM}}(\R;\Z)$.
Wartości grup zostaną uzasadnione przez dualność poniżej.
\end{przyklad}

\begin{przyklad}[Dlaczego jednopunktowe uzwarcenie bywa mylące?]
Dla nieskończonej przestrzeni dyskretnej $X=\N$
mamy $H_0^{\mathrm{BM}}(X;\Z)=\prod_{m\in\N}\Z$,
ponieważ lokalnie skończona suma punktów może mieć
dowolnie wiele niezerowych współczynników, a nie
ma nietrywialnych granic w stopniu zero.
Jednocześnie $H_0(X;\Z)=\bigoplus_{m\in\N}\Z$.
Nie wolno więc bez dodatkowych warunków utożsamiać
lokalnie skończonej homologii singularnej z
$H_*(X^+,\{\infty\};\Z)$: dla uzwarcenia
jednopunktowego przeliczalnej przestrzeni dyskretnej
ostatnia grupa w stopniu zero jest sumą prostą,
a nie produktem. Takie utożsamienie wymaga starannego
wyboru teorii homologii i hipotez o przestrzeni.
\end{przyklad}

\begin{przyklad}[Właściwy obraz może unicestwić klasę]
\label{prz:BM-proper-inclusion-line}
Włożenie $i:\R\hookrightarrow\R^2$,
$i(t)=(t,0)$, jest \emph{właściwe}:
przeciwobraz zwartego $K\subset\R^2$
jest zbiorem $K\cap(\R\times\{0\})$
identyfikowanym z zwartym podzbiorem prostej.
Indukuje więc dobrze określony homomorfizm
\[
 i_*:H_1^{\mathrm{BM}}(\R;\Z)\longrightarrow
      H_1^{\mathrm{BM}}(\R^2;\Z).
\]
Grupa po lewej jest izomorficzna z $\Z$,
natomiast grupa po prawej jest zerowa,
bo według dualności
$H_1^{\mathrm{BM}}(\R^2;\Z)
\cong H^1(\R^2;\Z)=0$.
A zatem $i_*=0$, mimo że $i$ jest
właściwym \emph{zanurzeniem}.

Geometrycznie generator po lewej jest
zorientowaną całą prostą; po włożeniu
w $\R^2$ staje się brzegiem
lokalnie skończonego łańcucha
triangulującego półpłaszczyznę
$\{(x,y)\in\R^2:y\ge0\}$.
Każdy zwarty zbiór przecina tylko skończenie
wiele trójkątów triangulacji lokalnie skończonej,
a składniki brzegów wewnętrznych znoszą się parami.
To pokazuje, dlaczego właściwość
gwarantuje istnienie $i_*$,
ale nie gwarantuje jego różnowartościowości.
\end{przyklad}

\section{Dualność dla rozmaitości niezwartych}
\begin{definicja}[Lokalnie skończona klasa fundamentalna]
Niech $M$ będzie spójną zorientowaną gładką
$n$-rozmaitością bez brzegu. Triangulacja lokalnie
skończona, zorientowana zgodnie z orientacją $M$,
wyznacza lokalnie skończony $n$-cykl złożony z
orientowanych sympleksów najwyższego wymiaru.
Jego klasa $[M]_{\mathrm{BM}}\in H_n^{\mathrm{BM}}(M;\Z)$
nazywa się klasą fundamentalną Borel--Moore'a.
Jej lokalny obraz w $H_n(M,M\setminus\{x\};\Z)$
jest lokalnym generatorem zgodnym z orientacją.
\end{definicja}

\begin{twierdzenie}[Dualność Poincarégo w wersji niezwartej]
\label{thm:noncompact-PD-v11}
Niech $M$ będzie gładką, drugoprzeliczalną,
zorientowaną $n$-rozmaitością bez brzegu,
a $G$ dowolną grupą abelową. Iloczyn z klasą
$[M]_{\mathrm{BM}}$ indukuje naturalne izomorfizmy
\[
H_c^k(M;G)\cong H_{n-k}(M;G),\qquad
H^k(M;G)\cong H_{n-k}^{\mathrm{BM}}(M;G)
\]
dla każdego $k\in\Z$. W przypadku rozmaitości
nieorientowalnej w grupach homologii po prawej
występuje system orientacji $\mathcal O_M\otimes_{\Z}G$.
\end{twierdzenie}
\begin{proof}
Przedstawimy model komórkowy realizujący oba izomorfizmy.
Gładka rozmaitość drugoprzeliczalna ma lokalnie skończoną
triangulację $K$. Po odpowiednim podziale barycentrycznym
każdemu zorientowanemu $k$-sympleksowi $\sigma$ odpowiada
blok dualny (komórka z podziału barycentrycznego) $D(\sigma)$ wymiaru $n-k$, złożona z
sympleksów barycentrycznych łańcuchów
$\sigma=\sigma_k<\sigma_{k+1}<\cdots<\sigma_n$.
Lokalna skończoność $K$ sprawia, że $D(\sigma)$ jest
skończonym łańcuchem i że rodzina komórek dualnych
jest lokalnie skończona.

Wybór orientacji $M$ wraz z orientacją $\sigma$
ustala orientację każdego bloku $D(\sigma)$.
W wymiarze zerowym blok dualny odpowiada
$n$-sympleksowi; dla lokalnie skończonej triangulacji
bloki są skończone, choć może ich być nieskończenie wiele.
Reguła orientacji brzegu implikuje,
że współczynnik $D(\tau)$ w $\partial D(\sigma)$
jest, z globalnie ustaloną konwencją znaków,
współczynnikiem $\sigma$ w $\delta\tau$.
W ten sposób operator przypisujący
$\sigma^*\mapsto D(\sigma)$ (gdzie $\sigma^*$ jest
kołańcuchem bazowym) daje izomorfizm kompleksów
z odpowiednim przesunięciem stopni między
komórkowymi kołańcuchami $K$ i dualnymi łańcuchami.

Zwykłe kołańcuchy komórkowe są \emph{dowolnymi}
rodzinami współczynników, tj.
$C^k(K;G)=\prod_{\sigma\in K^{(k)}}G$.
Odpowiadają one lokalnie skończonym łańcuchom
komórek dualnych. Kołańcuchy o nośniku zwartym
mają natomiast tylko skończenie wiele niezerowych
współczynników, czyli tworzą
$C_c^k(K;G)=\bigoplus_{\sigma\in K^{(k)}}G$;
odpowiadają zwykłym skończonym łańcuchom dualnym.
Komórkowo-singularne twierdzenie porównawcze,
stosowane do lokalnie skończonej triangulacji
i do łańcuchów lokalnie skończonych, pozwala
przejść od tych dwóch izomorfizmów kompleksów
do grup wymienionych w tezie. Konstrukcja jest
lokalna i odpowiada iloczynowi z klasą fundamentalną.

W przypadku nieorientowalnym orientacje komórek dualnych
nie sklejają się przy przejściu między mapami.
Ich zmiany znaków są dokładnie funkcjami przejścia
lokalnego systemu $\mathcal O_M$; ten system
musi wystąpić po stronie homologii.
\end{proof}

\begin{uwaga}[Zakres dowodu komórkowego]
Użyliśmy dwóch standardowych wyników topologii
rozmaitości: triangulowalności \emph{gładkich}
rozmaitości i zgodności homologii lokalnie
skończonej komórkowej z singularną. Nie rozszerzamy
tego argumentu bez dodatkowych hipotez na dowolne
niezwarte przestrzenie topologiczne ani na
niektóre nietriangulowalne rozmaitości topologiczne.
\end{uwaga}

\section{Przykłady i obliczenia grup}
\begin{przyklad}[Przestrzeń euklidesowa]
Dla $n\ge1$ znamy $H^0(\R^n;\Z)=\Z$ i
$H^k(\R^n;\Z)=0$ dla $k>0$.
Z dualności otrzymujemy
\[
H_k^{\mathrm{BM}}(\R^n;\Z)=
\begin{cases}\Z,&k=n,\\0,&k\ne n,\end{cases}
\qquad
H_c^k(\R^n;\Z)=
\begin{cases}\Z,&k=n,\\0,&k\ne n.\end{cases}
\]
Odtwarzamy w ten sposób rachunek form zwartych
z Twierdzenia~\ref{thm:compact-Rn-v11},
ale teraz również z całkowitymi współczynnikami.
\end{przyklad}

\begin{przyklad}[Otwarty cylinder]
Niech $M=S^1\times\R$, z orientacją produktową.
Ma typ homotopijny $S^1$, zatem $H_0(M)=H_1(M)=\Z$,
a pozostałe grupy homologii są zerowe.
Z dualności w wymiarze $2$ otrzymujemy
$H_c^0(M)=0$, $H_c^1(M)=\Z$, $H_c^2(M)=\Z$ oraz
$H_0^{\mathrm{BM}}(M)=0$, $H_1^{\mathrm{BM}}(M)=\Z$,
$H_2^{\mathrm{BM}}(M)=\Z$.
W wersji de Rhama generator stopnia $1$ reprezentuje
$\rho(t)\,\dd t$, a stopnia $2$ forma
$\dd\theta\wedge\rho(t)\,\dd t$, gdzie
$\int_{t\in\R}\rho(t)\,\dd t=1$ i $\dd\theta$ oznacza
standardową globalną formę kątową na okręgu.
\end{przyklad}

\begin{przyklad}[Nieskończenie wiele prostych: suma prosta kontra produkt]
\label{prz:BM-countable-lines}
Rozważmy rozłączną sumę topologiczną
\[
 X=\coprod_{j\in\N}\R_j,
\]
gdzie każda składowa $\R_j$ jest kopią prostej
z jej zwykłą topologią i orientacją.
Jest to \emph{niezwarta, drugoprzeliczalna
rozmaitość jednowymiarowa bez brzegu}, choć niespójna.

Każda składowa $\R_j$ jest otwarta i domknięta w $X$.
Jeśli $K\subset X$ jest zwarty, pokrycie
$K\subseteq\bigcup_{j\in\N}\R_j$
ma podpokrycie skończone; wobec tego
$K$ przecina \emph{tylko skończenie wiele składowych}.
Każdy sympleks singularny ma obraz spójny,
więc leży w jednej z $\R_j$.
Stąd, z uwzględnieniem nośników,
\[
 C_k(X;\Z)\cong\bigoplus_{j\in\N}C_k(\R_j;\Z),
 \qquad
 C_k^{\mathrm{lf}}(X;\Z)\cong
     \prod_{j\in\N}C_k^{\mathrm{lf}}(\R_j;\Z).
\]
Pierwsza równość wynika ze skończoności
zwykłych łańcuchów, a druga z faktu,
że lokalnie skończoną rodzinę łańcuchów
można wybrać \emph{niezależnie w każdej składowej}:
warunek na zwartych podzbiorach $X$
nie narzuca ograniczenia liczby składowych
z niezerowym współczynnikiem.

Różniczka działa współrzędnie.
Jądra i obrazy homomorfizmów iloczynowych
są iloczynami odpowiednich jąder i obrazów
(bo w każdej współrzędnej możemy osobno
wybrać przeciwobraz). Dlatego
\[
 H_1^{\mathrm{BM}}(X;\Z)
 \cong\prod_{j\in\N}
       H_1^{\mathrm{BM}}(\R_j;\Z)
 \cong\prod_{j\in\N}\Z.
\]
Z drugiej strony każda klasa o nośniku zwartym
jest reprezentowana w skończonej liczbie
składowych, a ekscyzja i definicja przez granicę
prostą dają
\[
 H_c^1(X;\Z)
 \cong\bigoplus_{j\in\N}H_c^1(\R_j;\Z)
 \cong\bigoplus_{j\in\N}\Z.
\]
Również zwykłe $H^0(X;\Z)\cong\prod_j\Z$,
bo funkcję lokalnie stałą wybieramy oddzielnie
w każdej składowej, podczas gdy
$H_0(X;\Z)\cong\bigoplus_j\Z$.

Obie wersje dualności
z Twierdzenia~\ref{thm:noncompact-PD-v11}
odtwarzają dokładnie te wyniki:
\[
 H_c^1(X;\Z)\cong H_0(X;\Z),
 \qquad
 H^0(X;\Z)\cong H_1^{\mathrm{BM}}(X;\Z).
\]
Grup tych \emph{nie wolno utożsamiać}:
suma prosta składa się z ciągów o skończonej
liczbie niezerowych wyrazów, a produkt
dopuszcza dowolne ciągi całkowite.
Na przykład lokalnie skończony cykl złożony
ze zorientowanej całej prostej w każdej
składowej odpowiada elementowi
$(1,1,1,\ldots)$ produktu, który
nie należy do sumy prostej.
\end{przyklad}

\begin{figure}[!htbp]
\centering
\begin{tikzpicture}[>=Stealth,font=\small]
  \foreach \y/\j in {2.5/1,1.5/2,.5/3} {
    \draw[manifoldblue,thick] (-3.6,\y)--(3.6,\y);
    \foreach \x in {-3,-2,-1,0,1,2} {
      \draw[->,manifoldblue,thick]
          (\x,\y+.1)--(\x+.70,\y+.1);
    }
    \node[right] at (3.8,\y) {$\R_{\j}$};
  }
  \node[align=center] at (0,-.45)
       {$\vdots$\\[3pt] $\R_j\quad(j\in\N)$};
  \draw[manifoldred,very thick,dashed,rounded corners]
       (-1.4,1.20) rectangle (1.4,2.82);
  \node[above,manifoldred] at (0,2.88) {$K$};
\end{tikzpicture}
\caption{Przeliczalna rozłączna suma prostych.
Przykładowy zwarty podzbiór $K$ przecina tylko
skończenie wiele składowych. Lokalnie skończony
cykl może mieć jednak niezerowy współczynnik
przy \emph{każdej} z nieskończenie wielu prostych.}
\label{fig:BM-countable-lines}
\end{figure}

\begin{przyklad}[Wnętrze zwartej rozmaitości z brzegiem]
Niech $W$ będzie zwartą gładką $n$-rozmaitością
z brzegiem i niech $M=\operatorname{int}W$.
Kołnierz brzegu i ekscyzja, w połączeniu z
porównaniem komórkowym, dają
$H_c^k(M;G)\cong H^k(W,\partial W;G)$ oraz
$H_k^{\mathrm{BM}}(M;G)\cong H_k(W,\partial W;G)$.
Dla $W=D^n$ uzyskujemy
$H_n^{\mathrm{BM}}(\R^n;\Z)=\Z$,
choć zwykła homologia $H_n(\R^n;\Z)$ jest zerowa.
\end{przyklad}

\section{Związek z dualnością Poincarégo--Lefschetza}
Jeśli $M=\operatorname{int}W$ dla zwartej rozmaitości
z brzegiem, to izomorfizmy
$H_c^k(M)\cong H^k(W,\partial W)$ oraz
$H_{n-k}(M)\cong H_{n-k}(W)$
zamieniają pierwszą dualność z
Twierdzenia~\ref{thm:noncompact-PD-v11}
w dualność Poincarégo--Lefschetza dla pary
$(W,\partial W)$. Druga dualność odpowiada
$H^k(W)\cong H_{n-k}(W,\partial W)$.
W ten sposób oba równania nie są osobnymi
przypadkowymi analogiami: są dwiema wersjami
tego samego mechanizmu geometrycznego.

\section{Wskazówki przed połączeniem z geometrią algebraiczną}
Dla gładkiej zespolonej rozmaitości algebraicznej
wymiaru zespolonego $d$ występuje naturalna orientacja
rzeczywistej $2d$-rozmaitości. Grupy
$H_{2d}^{\mathrm{BM}}$ są właściwym miejscem dla klas
fundamentalnych niezwartych podrozmaitości, a
$H_c^*(\cdot)$ jest potrzebna do dualności i map
właściwych. W następnym wspólnym skrypcie można
połączyć te konstrukcje z cyklami algebraicznymi,
przecięciami i kohomologią snopów; utożsamienie
cykli algebraicznych z wszystkimi klasami topologicznymi
nie jest jednak prawdziwe bez dodatkowych warunków.


